\section{国际投资法律制度}
\begin{enumerate}
  \item 多边投资担保机构公约。鼓励\zd{资本向发展中国家流动}。
        \begin{kk}{多边投资担保机构/MIGA}
          是一个国际\zd{保险机构}，独立的国际法主体。提供\zd{政治风险、国家风险}担保。
          \begin{enumerate}
            \item 保险资金来源于各国政府的认缴。
            \item 投保人为\zd{私人投资者}（自然人、法人，国家不行）。投到发展中会员国（称为\zd{合格东道国}）。
            \item 主要 4 种风险。
                  \begin{enumerate}
                    \item 政治引发的\zd{货币汇兑}风险。不包括正常的商业风险。\al{货币贬值 \cuo}。
                    \item \zd{征收}风险。\yi 没收。\er 丧失本应获得的经济利益。
                    \item \zd{战争内乱}。\al{罢工 \cuo}。
                    \item \zd{政府违约且拒绝司法}。
                  \end{enumerate}
            \item 先投保后投资。\zd{理赔要求用尽当地救济}（行政、司法）。
          \end{enumerate}
        \end{kk}

  \item 解决国家和其他国民之间投资争端公约。
        \begin{kk}{国际投资争端解决中心/ICSID。本质是一个仲裁机构}
          \begin{enumerate}
            \item 国家和公民之间。
            \item 只受理国际投资\zd{法律争端}。
            \item \zd{不以用尽当地救济为条件}，双方出具\zd{书面同意}文件即可。
            \item 排他性。一旦确立了中心的管辖权，就排除其他争端解决方式。
            \item 一裁终局，有约束力。
          \end{enumerate}
        \end{kk}
\end{enumerate}
